\section{Kết luận và hướng phát triển}
\label{sec:conclusion}

\subsection{Kết luận chung}

Báo cáo đã trình bày một pipeline hoàn chỉnh cho bài toán phân loại âm thanh đô
thị trên UrbanSound8K bằng học sâu, đi từ xác định bài toán, phân tích dữ liệu,
thiết kế tiền xử lý, xây dựng \texttt{Dataset}/\texttt{DataLoader}, huấn luyện mô
hình cho tới đánh giá và phân tích kết quả.

Một bài học xuyên suốt là vai trò của bước tiền xử lý. Dữ liệu gốc không đồng
nhất rõ rệt về tần số lấy mẫu, số kênh và độ dài; nếu bỏ qua, mô hình có nguy cơ
học nhầm định dạng thay vì nội dung âm học. Các quyết định rút ra từ EDA -- chuyển
mono, resample về 22{,}05\,kHz, chuẩn hóa 4{,}0 giây và trích xuất log-Mel
spectrogram -- đã tạo nền tảng ổn định để bốn cấu hình được so sánh một cách công
bằng.

\subsection{Những kết quả nổi bật}

Đề tài đã thực hiện và đánh giá bốn cấu hình: CNN baseline, ResNet18, ResNet18 +
augmentation, và ensemble của hai mô hình ResNet. Các phát hiện chính:

\begin{itemize}[leftmargin=1.4em,itemsep=2pt]
  \item \textbf{ResNet18 vượt CNN baseline} ở cả Accuracy ($0{,}7434$ so với
        $0{,}6999$) lẫn Macro-F1 ($0{,}7625$ so với $0{,}7113$). Trên bài toán
        này, độ sâu và kết nối phần dư thực sự có ích.
  \item \textbf{Augmentation là một sự đánh đổi, không phải món quà miễn phí.}
        Khi dùng đơn lẻ, nó hạ nhẹ hiệu năng tổng thể của ResNet18, nhưng tái
        phân bố năng lực: cải thiện các lớp trầm như \texttt{air\_conditioner},
        đồng thời làm giảm độ sắc ở nhóm cơ khí, rõ nhất là \texttt{jackhammer}.
  \item \textbf{Ensemble cho kết quả tốt nhất} -- Accuracy $0{,}7572$, Macro-F1
        $0{,}7759$ -- và là cấu hình duy nhất vượt mốc $0{,}75$. Nó thành công
        nhờ ghép hai mô hình có hồ sơ lỗi khác nhau, chứ không nhờ ghép hai mô
        hình cùng mạnh.
\end{itemize}

\subsection{Ý nghĩa học thuật và thực tiễn}

Về mặt học thuật, đề tài cho thấy hiệu năng trong phân loại âm thanh phụ thuộc
vào nhiều yếu tố hơn là chỉ kiến trúc: chất lượng tiền xử lý, mức độ augmentation
và cách kết hợp mô hình đều để lại dấu ấn rõ trên kết quả. Trường hợp augmentation
là một ví dụ đáng nhớ -- một kỹ thuật thường được mặc định là ``có lợi'' lại có
thể gây hại nếu không xét tới bản chất từng lớp dữ liệu.

Về mặt thực tiễn, project để lại một hệ thống có cấu trúc rõ ràng và tái lập tốt.
Các thành phần tiền xử lý, \texttt{Dataset}, huấn luyện, đánh giá và logging W\&B
đều tách biệt, có thể tái sử dụng hoặc mở rộng cho các nghiên cứu tiếp theo.

\subsection{Hướng phát triển}

Nếu tiếp tục, đề tài có thể được mở rộng theo các hướng sau:

\begin{itemize}[leftmargin=1.4em,itemsep=2pt]
  \item \textbf{Kiểm định chéo 10-fold} theo đúng giao thức chuẩn của
        UrbanSound8K, để có con số trung bình kèm độ dao động thay vì kết quả
        trên một fold duy nhất.
  \item \textbf{Đặc trưng giàu hơn} cho nhóm cơ khí -- chẳng hạn ghép thêm các
        kênh delta và delta-delta -- nhằm tách \texttt{air\_conditioner},
        \texttt{engine\_idling}, \texttt{drilling} và \texttt{jackhammer}.
  \item \textbf{Mô hình tiền huấn luyện} quy mô lớn như PANNs, được tinh chỉnh
        trên UrbanSound8K, để nâng trần hiệu năng.
  \item \textbf{Khảo sát augmentation có hệ thống} nhằm tìm cấu hình giữ được
        lợi ích cho lớp trầm mà không phá đặc trưng lặp của nhóm cơ khí.
\end{itemize}

\subsection{Thông tin nộp bài}

Mã nguồn, notebook và toàn bộ nhật ký thí nghiệm được công khai tại:

\begin{itemize}[leftmargin=1.4em,itemsep=1pt]
  \item Mã nguồn (GitHub): \\
        \url{https://github.com/Zackerville/UrbanSoundClassification_DeepLearning_8K}
  \item Nhật ký thí nghiệm (Weights \& Biases): \\
        \url{https://wandb.ai/tien-ngozack2004-ho-chi-minh-city-university-of-technology/urban-sound-classification}
\end{itemize}
